\documentclass[10pt, a4paper]{article}

% Parametri che modificano il file main.tex
% Le uniche parti da cambiare su main.tex sono:
% - tabella versioni
% - testo file

\def\titolo{Analisi dei \\ \ requisiti} % \\ per andare a capo, \ per spazio
\def\titoloHeader{Analisi dei requisiti}
\def\data{v2.0.0(0)}

% \def\listaComponenti{
% Bresolin G.,
% Campese M.,
% Ciriolo I.,
% Dugo A.,
% Feltrin E.,
% Michelon R.,
% Orlandi G.
% }

\usepackage{style}
\usepackage{headerfooter}
\usepackage{caption}
\usepackage{hyperref}
\usepackage{svg}
\usepackage{comment}
\usepackage{setspace}
\definecolor{primarycolor}{RGB}{248, 182, 143}

\title{\titolo}
\author{SWEetCode}

\begin{document}

% PRIMA PAGINA
\begin{titlepage}
    \thispagestyle{empty}
    \begin{tikzpicture}[remember picture, overlay]
        % TRIANGOLI
        \draw[fill=secondarycolor, secondarycolor] (current page.north west) -- (current page.south west) -- (8.8, -28);
        \draw[fill=primarycolor, primarycolor] (-3, 5) -- (4, -13.6) -- (11, 5);

        % LOGO
        \node [xshift=-5cm, yshift=25cm] (logo) at (current page.south east) {\includegraphics[width=6.5cm]{img/logo.png}};

        % SWEETCODE - DATE
        \node [anchor=north east, align=right, xshift=-1.2cm, yshift=20.5cm, text=black] (sweetcode) at (current page.south east) {\fontsize{32pt}{36pt}\selectfont SWEetCode};
        \draw[line width=4pt, lightcol] ([xshift=-3cm, yshift=-0.37cm]sweetcode.south west) -- ([yshift=-0.37cm]sweetcode.south east);
        \node [anchor=north east, align=right, xshift=-1.2cm, yshift=18.7cm, text=black] (date) at (current page.south east){\fontsize{24pt}{24pt} \selectfont Verbale \tipoVerb};

        % NOME FILE
        \node [anchor=north east, text width=15cm, align=right, xshift=-1.2cm, yshift=17cm, text=black] (titolo) at (current page.south east){\fontsize{48pt}{48pt}\textbf{\data}};

        % BOX DATI PARTECIPANTI
        \ifthenelse{\equal{\tipoVerb}{Esterno}}{
            \node[anchor=north east, xshift=-1.2cm, yshift=14.5cm, minimum width=8cm] (box) at (current page.south east){};
        }{
            \node[anchor=north east, xshift=-1.2cm, yshift=12.5cm, minimum width=8cm] (box) at (current page.south east){};
        }

        % RESPONSABILE
        \node[anchor=north west, align=left] (dati2) at (box.north west) {\fontsize{15pt}{15pt}\selectfont \textbf{Responsabile}};
        \draw[line width=4pt, lightcol] (dati2.south west) -- ([xshift=8cm]dati2.south west);
        \node[anchor=north west, align=left] (dati21) at (dati2.south west){\fontsize{13pt}{13pt}\selectfont \nomeResp};

        % VERIFICATORE
        \node[anchor=north west, yshift=-1cm, align=left] (dati3) at (dati21.north west) {\fontsize{15pt}{15pt}\selectfont \textbf{Verificatore}};
        \draw[line width=4pt, lightcol] (dati3.south west) -- ([xshift=8cm]dati3.south west);
        \node[anchor=north west, align=left] (dati31) at (dati3.south west){\fontsize{13pt}{13pt}\selectfont \nomeVer};

        % SEGRETARIO DI RIUNIONE
        \node[anchor=north west, yshift=-1cm, align=left] (dati4) at (dati31.north west) {\fontsize{15pt}{15pt}\selectfont \textbf{Segretario di Riunione}};
        \draw[line width=4pt, lightcol] (dati4.south west) -- ([xshift=8cm]dati4.south west);
        \node[anchor=north west, align=left] (dati41) at (dati4.south west){\fontsize{13pt}{13pt}\selectfont \nomeSegr};
        
        % UNIPD - SWE
        \node [xshift=4.4cm, yshift=2.3cm, draw, secondarycolor, text=white] (uni) at (current page.south west) {\fontsize{20pt}{20pt} \selectfont Università di Padova};
        \node [xshift=0.65cm, yshift=0.7cm, draw, secondarycolor, text=white, below=of uni] (corso) {\fontsize{20pt}{20pt}\selectfont Ingegneria del Software};

        % FIRMA AZIENDA
        \ifthenelse{\equal{\tipoVerb}{Esterno}}{
            \draw[line width=4pt, lightcol] ([xshift=-1.2cm, yshift=7cm]current page.south east) -- ([xshift=-8cm, yshift=7cm]current page.south east);
            \node [xshift=-4.8cm, yshift=8cm] (logo) at (current page.south east) {\includegraphics[width=6cm]{img/firme/\firmaAzienda}};
            \node[anchor=north west, xshift=12.9cm, yshift=6.65cm, align=left] at (current page.south west)
            {\fontsize{13pt}{13pt}\selectfont L'azienda: \nomeAzienda};
        }

        % FIRMA CODING COWBOYS
        \draw[line width=4pt, lightcol] ([xshift=-1.2cm, yshift=4.35cm]current page.south east) -- ([xshift=-8cm, yshift=4.35cm]current page.south east);
        \node [xshift=-4.8cm, yshift=5.5cm] (logo) at (current page.south east) {\includegraphics[width=6cm]{img/firme/menon.png}};
        \node[anchor=north west, xshift=12.9cm, yshift=4cm, align=left] at (current page.south west)
        {\fontsize{13pt}{13pt}\selectfont Coding Cowboys: Menon G.};

        % FIRMA
        \draw[line width=4pt, lightcol] ([xshift=-1.2cm, yshift=1.8cm]current page.south east) -- ([xshift=-8cm, yshift=1.8cm]current page.south east);
        \node [xshift=-4.8cm, yshift=2.45cm] (logo) at (current page.south east) {\includegraphics[width=6cm]{img/firme/\firmaResp}};
        \node[anchor=north west, xshift=12.9cm, yshift=1.45cm, align=left] at (current page.south west)
        {\fontsize{13pt}{13pt}\selectfont Il Responsabile: \nomeResp};
        
    \end{tikzpicture}
\end{titlepage}

% REGISTRO DELLE VERSIONI
%Registro in ordine dalla più recente alla meno recente!
{\renewcommand{\arraystretch}{1.5}
\section*{Registro delle versioni}

\begin{xltabular}{\textwidth}{c|c|c|c|X}
\label{tab:long}

\textbf{Versione} & \textbf{Data} & \quantities{\textbf{Responsabile di}\\\textbf{stesura}}& \textbf{Revisore} & \quantities{\textbf{Dettaglio e}\\\textbf{motivazioni}} \\
\endfirsthead

\textbf{Versione} & \textbf{Data} & \quantities{\textbf{Responsabile di}\\\textbf{stesura}}& \textbf{Revisore} & \quantities{\textbf{Dettaglio e}\\\textbf{motivazioni}} \\
\endhead

\multicolumn{5}{r}{{Continua nella pagina successiva}} \\
\endfoot

\endlastfoot

\hline
v1.10.0(7) & $2023-12-29$ & \quantities{Michelon R.} & Dugo A. & Aggiunti diagrammi use case e di attività.\\
\hline
v1.10.0(6) & $2023-12-29$ & \quantities{Michelon R.} & Dugo A. & Aggiornamento use case e requisiti.\\
\hline
v1.0.0(14) & $2023-12-06$ & \quantities{Bresolin G.} & Orlandi G. & Aggiornamento UC e requisiti presenti. Introduzione diagrammi dei casi d'uso. Aggiunta di nuovi UC e nuovi requisiti.\\
\hline
v1.0.0(5) & $2023-11-23$ & \quantities{Ciriolo I.} & Campese M. & Introduzione, Descrizione, Studio dei primi Casi d'uso, Studio dei primi requisiti.\\

    
\end{xltabular}}
%\newpage
\newpage
%todo modificare questa section per avere sempre un titolo "grande" ma senza numero
\section*{Candidatura per Requirements and Tecnology Baseline}
Il gruppo SWEetCode è lieto di presentare la propria candidatura per la revisione di avanzamento Requirements and Tecnology Baseline (RTB), nell'impegno di consegnare il progetto da Voi commissionato, denominato \textbf{Knowledge Management AI} (C1), proposto dall'azienda \textbf{AzzurroDigitale}.\\

Per una panoramica delle attività e una consultazione semplificata dei documenti riguardanti l'RTB, il team di SWEetCode estende un invito cordiale a visitare il sito:
\begin{center}
\href{https://sweetcode-team.github.io}{\textcolor{blue}{sweetcode-team.github.io}}
\end{center}


In tale spazio si potranno trovare:
\begin{itemize}
    \item (Analisi dei requisiti v2.0.0(0));
    \item (Glossario v2.0.0(0));
    \item (Norme di progetto v2.0.0(0));
    \item (Piano di progetto v2.0.0(0));
    \item (Piano di qualifica v2.0.0(0));
    \item I verbali delle riunioni interne tenute i giorni:
        \begin{itemize}
            \item 2023-11-07;
            \item 2023-11-17;
            \item 2023-11-22;
            \item 2023-11-27;
            \item 2023-12-06;
            \item 2023-12-27;
            \item 2024-01-15;
            \item 2024-02-03;
        \end{itemize}
   \item I verbali delle riunioni esterne tenute i giorni:
        \begin{itemize}
            \item 2023-11-09: AzzurroDigitale;
            \item 2023-11-15: AzzurroDigitale;
            \item 2023-11-22: AzzurroDigitale;
            \item 2023-11-24: AzzurroDigitale;
            \item 2023-11-29: AzzurroDigitale;
            \item 2023-12-06: AzzurroDigitale;
            \item 2024-01-24: AzzurroDigitale;
        \end{itemize}    
\end{itemize}

Per consultare il repository del progetto, il team SWEetCode estende un invito cordiale a visitare il link:\\
\begin{center}
\href{https://github.com/sweetcode-team/Knowledge_Management_AI}{\textcolor{blue}{github.com/sweetcode-team/Knowledge\_Management\_AI}}
\end{center}

\newpage
\subsection*{Aggiornamento Impegni}
Il team procede nello sviluppo del progetto verso la Product Baseline confermando la data di consegna prevista per il 2024-04-08, ma diminuendo il preventivo da 12845,00€ a 12715,00€,  a seguito del cambiamento dei quantitativi orari assegnati ai ruoli di responsabile, amministratore, analista, progettista e programmatore.
Le tabelle orarie e le motivazioni dei cambiamenti sono riportate nel documento (Piano di progetto v2.0.0(0), \S Stima dei costi di realizzazione).

{\renewcommand{\arraystretch}{1.5}
\begin{table}[H]
\centering
\begin{tabularx}{0.54\textwidth}{c|c|c}
\rowcolor{primarycolor}
    \textbf{Revisione} & \textbf{Data}& \textbf{Costi}\\
\hline
    \textbf{Product Baseline} & 2024-04-08 & 12715,00€\\
\end{tabularx}
\end{table}

\vspace{1em}
Il team è composto dai seguenti membri:
\\
\\
% TABELLA (tenere anche arraystretch che aumenta il padding verticale)
{\renewcommand{\arraystretch}{1.5}
\begin{tabularx}{\textwidth}{X|X}
\textbf{Cognome e Nome} & \textbf{Matricola n°} \\
\hline
Bresolin Gianluca & 2034316 \\
\hline
Campese Martina & 1122439 \\
\hline
Ciriolo Irene & 2043682 \\
\hline
Dugo Alberto & 2042382 \\
\hline
Feltrin Emanuele & 2034314 \\
\hline
Michelon Riccardo & 2042341 \\
\hline
Orlandi Giacomo & 1122253 \\
\end{tabularx}}\\

Cordiali saluti,\\

Il gruppo SWEetCode

\end{document}